{\Large\bfseries Midterm Exam\par}
\bigskip

\fullname

The exam contains 8 questions, worth 2 to 4 points each as marked. The total is 22 points and the score is out of 20 (so you still get full credit even if you lose up to 2 points).

60 minutes. One handwritten A6 (or 4x6 index card) cheat sheet allowed, no printed material.
Remember that your electronic devices should be silent and put away in your bag at all times.

\bigskip
\begin{center}\rule{0.4\linewidth}{0.4pt}\end{center}
\bigskip

\begin{question}
(2 points) Consider the following class, and the code fragment below it, which runs inside a method.

\begin{lstlisting}
class Node {
    int data;
    Node next;
    Node(int data, Node next) { this.data = data; this.next = next; }
}
\end{lstlisting}

\begin{lstlisting}
Node head = null;
for (int i = 1; i <= 4; i++)
    head = new Node(i, head);
Node keep = head.next.next;
head = head.next;
head.next = null;
\end{lstlisting}

(a) Right after the last line is executed, which \lstinline|Node| objects can be garbage collected? Identify them by their \lstinline|data|.

Answer: \underline{\hspace{2in}} \ifanswers (only the node with data 4: \lstinline|head| refers to 3, \lstinline|keep| to 2, which refers to 1, and nothing refers to 4)\fi

\bigskip
(b) Right after the last line is executed, mark where each of the following is stored.

\begin{center}
\begin{tabular}{@{}p{3.6in}cc@{}}
\toprule
 & Stack & Heap \\
\midrule
the variable \lstinline|keep| & \cb{1} & \cb{0} \\
the \lstinline|Node| object whose \lstinline|data| is 3 & \cb{0} & \cb{1} \\
the field \lstinline|next| of the \lstinline|Node| whose \lstinline|data| is 2 & \cb{0} & \cb{1} \\
the \lstinline|int| value 1, stored in the field \lstinline|data| of a \lstinline|Node| & \cb{0} & \cb{1} \\
\bottomrule
\end{tabular}
\end{center}
\end{question}

\begin{question}
(2 points) In each fragment below, \lstinline|x| holds $N$ elements before the fragment runs, and was declared as one of \lstinline|ArrayList<Integer> x|, \lstinline|LinkedList<Integer> x| or \lstinline|ArrayDeque<Integer> x|.
\lstinline|sum| is an \lstinline|int| variable.

\begin{minipage}[t]{0.48\linewidth}
(A)
\begin{lstlisting}
for (int i = 0; i < N; i++)
    x.push(i);
while (!x.isEmpty())
    sum += x.pop();
\end{lstlisting}
\end{minipage}\hfill
\begin{minipage}[t]{0.48\linewidth}
(B)
\begin{lstlisting}
for (int i = 0; i < x.size(); i++)
    sum += x.get(i);
\end{lstlisting}
\end{minipage}

For each fragment and each declaration, write ``does not compile'' if the fragment does not compile, and otherwise the worst-case complexity of the whole fragment, as a function of $N$.

\begin{center}
\renewcommand{\arraystretch}{2}
\begin{tabular}{@{}c|p{1.4in}|p{1.4in}|p{1.4in}@{}}
 & \lstinline|ArrayList| & \lstinline|LinkedList| & \lstinline|ArrayDeque| \\
\hline
(A) & \ifanswers does not compile\fi & \ifanswers $O(N)$\fi & \ifanswers $O(N)$\fi \\
\hline
(B) & \ifanswers $O(N)$\fi & \ifanswers $O(N^2)$\fi & \ifanswers does not compile\fi \\
\end{tabular}
\end{center}

\ifanswers (A: \lstinline|ArrayList| has no \lstinline|push|/\lstinline|pop|; on \lstinline|ArrayDeque| a single \lstinline|push| can cost $O(N)$, but $N$ of them cost $O(N)$ in total. B: \lstinline|ArrayDeque| has no \lstinline|get(i)|.)\fi
\end{question}

\begin{question}
(2 points) The method below computes the text typed on a keyboard, given the sequence of keys pressed.
The key \lstinline|<| is a backspace: it erases the last character of the text typed so far (if the text is empty, it does nothing).
For example, \lstinline|typedText("cat<<ow")| returns \lstinline|"cow"|, and \lstinline|typedText("a<<b")| returns \lstinline|"b"|.

\begin{lstlisting}
String typedText(String keys) {
    Deque<Character> text = new ArrayDeque<>();
    for (char c : keys.toCharArray()) {
        if (c == '<') {
            if (!text.isEmpty())
                text.______________;        // (1)
        } else {
            text.______________;            // (2)
        }
    }
    StringBuilder sb = new StringBuilder();
    while (!text.isEmpty())
        sb.append(text.______________);     // (3)
    return sb.toString();
}
\end{lstlisting}

Fill in the blanks with methods of \lstinline|Deque|.

(1): \lstinline|text.|\underline{\hspace{2in}} \ifanswers (\lstinline|removeLast()|)\fi

\medskip
(2): \lstinline|text.|\underline{\hspace{2in}} \ifanswers (\lstinline|addLast(c)|)\fi

\medskip
(3): \lstinline|text.|\underline{\hspace{2in}} \ifanswers (\lstinline|removeFirst()|)\fi

\ifanswers (Also correct, mirrored: \lstinline|removeFirst()|, \lstinline|addFirst(c)| (or \lstinline|push(c)|), \lstinline|removeLast()|. Blanks (1) and (2) must work on the same end, (3) on the other one.)\fi
\end{question}

\begin{question}
(2 points) Consider the following class, and three comparators for it.

\begin{lstlisting}
class Student {
    String name;
    int section;
    int score;
}

class BySection   implements Comparator<Student> { ... }   // section, smallest first
class ByScoreDesc implements Comparator<Student> { ... }   // score, highest first
class ByName      implements Comparator<Student> { ... }   // name, alphabetical
\end{lstlisting}

The array \lstinline|roster| must be ordered by section; students in the same section by score, from highest to lowest; students with the same section and the same score by name.
Instead of writing a new comparator, we call \lstinline|sort(roster, c)| three times, once with each comparator above.
\lstinline|sort| uses one of the sorting algorithms seen in class.

(a) What property should \lstinline|sort| have for the three calls to produce the required ordering?

Answer: \underline{\hspace{1.5in}} \ifanswers (stable: each call must keep the order produced by the previous calls among students that it considers equal)\fi

\bigskip
(b) Name one of the algorithms seen in class that has that property.

Answer: \underline{\hspace{1.5in}} \ifanswers (bubble sort, insertion sort or merge sort)\fi

\bigskip
(c) Write the numbers 1 to 3 in the boxes to give the order in which the three calls must be executed.

\medskip
\begin{tabular}{@{}l@{\quad}l@{}}
\ansbox{3} & \lstinline|sort(roster, new BySection());| \\[6pt]
\ansbox{1} & \lstinline|sort(roster, new ByName());| \\[6pt]
\ansbox{2} & \lstinline|sort(roster, new ByScoreDesc());| \\
\end{tabular}
\end{question}

\begin{question}
(3 points) We say that an array of $N$ distinct elements is \textbf{$k$-sorted} if every element is at most $k$ positions away from its position in the sorted array.
For example, \lstinline|{2, 1, 3, 5, 4}| is 1-sorted, and \lstinline|{3, 1, 2, 6, 4, 5}| is 2-sorted.

(a) What is the worst-case complexity of insertion sort on a $k$-sorted array, as a function of $N$ and $k$? Select one, and justify your answer briefly.

\cb{0} $O(N)$ \qquad \cb{0} $O(k)$ \qquad \cb{0} $O(\min(N, k))$ \qquad \cb{0} $O(N + k)$ \qquad \cb{1} $O(Nk)$

\begin{answerspace}{1in}
\ifanswers
Inserting \lstinline|a[i]| moves it past the larger elements before it, and there are at most $k$ of those: the $i - k$ smallest elements belong at indices $0..i-k-1$, so they all start before index $i$, and they are all smaller than \lstinline|a[i]|. So each insertion costs $O(k)$.
(Accept ``each element moves at most $k$ positions'' as a justification.)
\fi
\end{answerspace}

\bigskip
(b) What is the worst-case complexity of merge sort on a $k$-sorted array?

Hint: merge sort is not adaptive.

Answer: \underline{\hspace{1in}} \ifanswers ($O(N \log N)$)\fi

\bigskip
(c) Let $N = 1{,}000{,}000$. Given your answers above, which algorithm would you choose for $k = 5$, and which for $k = 1000$? Assume that the constants and low-order terms hidden in the two complexities above are similar.

Hint: $\log_2 N \approx 20$.

\begin{center}
\begin{tabular}{@{}p{1.2in}cc@{}}
\toprule
 & insertion sort & merge sort \\
\midrule
$k = 5$ & \cb{1} & \cb{0} \\
$k = 1000$ & \cb{0} & \cb{1} \\
\bottomrule
\end{tabular}
\end{center}
\ifanswers ($k = 5$: about $5 \cdot 10^6$ vs. $2 \cdot 10^7$; $k = 1000$: about $10^9$ vs. $2 \cdot 10^7$)\fi
\end{question}

\begin{question}
(3 points) Consider an array-based queue implementation that stores its elements in \lstinline|data[front]|, \ldots, \lstinline|data[front + size - 1]|, \textbf{without} wrapping around (it does \textbf{not} use a circular array):
\begin{itemize}[leftmargin=2em,itemsep=1pt,topsep=2pt]
\item \lstinline|enqueue(e)| stores \lstinline|e| in \lstinline|data[front + size]|; if \lstinline|front + size == data.length|, it first doubles the capacity of the array;
\item \lstinline|dequeue()| returns \lstinline|data[front]| and increments \lstinline|front|.
\end{itemize}
To reuse the empty cells before \lstinline|front|, the queue sometimes \textbf{compacts}: it copies its elements to \lstinline|data[0]|, \ldots, \lstinline|data[size - 1]| and sets \lstinline|front = 0|, in $O($\lstinline|size|$)$ time.

(a) Consider the following three strategies for when to compact, and select the one that makes both \lstinline|enqueue| and \lstinline|dequeue| $O(1)$ amortized:
\begin{itemize}[leftmargin=2em,itemsep=1pt,topsep=2pt]
\item[\cb{0}] at the end of every \lstinline|dequeue()|;
\item[\cb{1}] at the end of a \lstinline|dequeue()|, if \lstinline|front >= size|;
\item[\cb{0}] at the start of an \lstinline|enqueue(e)|, if \lstinline|front > 0| and \lstinline|front + size == data.length| (i.e. compacting instead of doubling).
\end{itemize}

\ifanswers (A compaction moves \lstinline|size| elements, and since the previous one \lstinline|front| went from 0 to at least \lstinline|size|, so at least \lstinline|size| dequeues paid for it. Compacting at every dequeue costs $O(N)$ each time.)\fi

\bigskip
(b) The queue starts empty, with an array of capacity 8, and the following operations are performed, using the strategy you selected in (a):

\begin{minipage}[t]{0.35\linewidth}
\vspace{-\baselineskip}
\begin{lstlisting}
for (int i = 1; i <= 5; i++)
    q.enqueue(i);
q.dequeue();
q.dequeue();
q.dequeue();
q.enqueue(6);
q.dequeue();
\end{lstlisting}
\end{minipage}\hfill
\begin{minipage}[t]{0.6\linewidth}
Draw the final state of the array (every cell), and give the final values of \lstinline|front| and \lstinline|size|.

Hint: assume that the cells that do not hold an element of the queue contain \lstinline|null|.
\end{minipage}

\begin{answerspace}{1.2in}
\ifanswers
\begin{verbatim}
data (cells 0..7): [null, 5, 6, null, null, null, null, null]
front = 1, size = 2
(the 3rd dequeue makes front = 3 >= size = 2, so it compacts:
[4, 5, ...], front = 0; the last dequeue makes front = 1 < size = 2,
so it does not compact)
With the other strategies, for follow-through credit:
compact at every dequeue: [5, 6, null, null, null, null, null, null]
                          front = 0, size = 2
compact in enqueue:       [null, null, null, null, 5, 6, null, null]
                          front = 4, size = 2
\end{verbatim}
\fi
\end{answerspace}
\end{question}

\begin{question}
(4 points) Consider the following method, which uses a \lstinline|Deque| as a stack.

\begin{lstlisting}
void printIter(int n) {
    Deque<Integer> stack = new ArrayDeque<>();
    stack.push(n);
    while (!stack.isEmpty()) {
        int m = stack.pop();
        if (m <= 0) continue;
        System.out.print(m + " ");
        stack.push(m - 1);
        stack.push(m - 2);
    }
}
\end{lstlisting}

(a) What does \lstinline|printIter(4)| print?

Answer: \underline{\hspace{2.5in}} \ifanswers (\lstinline|4 2 1 3 1 2 1|)\fi

\bigskip
(b) Write a recursive method \lstinline|printRec(int n)| that prints exactly the same as \lstinline|printIter(n)|, for every \lstinline|n|, without using any loop or \lstinline|Deque|.

\begin{answerspace}{1.8in}
\ifanswers
\begin{lstlisting}
void printRec(int n) {
    if (n <= 0) return;
    System.out.print(n + " ");
    printRec(n - 2);
    printRec(n - 1);
}
\end{lstlisting}
The calls are in the opposite order of the pushes: \lstinline|m - 2| is pushed last, so it is popped first, and everything it pushes is popped before \lstinline|m - 1|.
With the calls in push order, \lstinline|printRec(4)| would print \lstinline|4 3 2 1 1 2 1|.
\fi
\end{answerspace}
\end{question}

\begin{question}
(4 points) A singly linked list has the following fields.

\begin{lstlisting}
public class MyLinkedList<E> {
    private static class Node<E> {
        E data;
        Node<E> next;
    }
    private Node<E> head;    // first node, null if the list is empty
    private Node<E> tail;    // last node, null if the list is empty
    private int size;
    ...
}
\end{lstlisting}

Write the method \lstinline|public void rotateRight(int k)|, where $k \ge 0$, which moves the last \lstinline|k| elements to the front of the list, keeping their order.
Rotating by \lstinline|size| gives back the same list, so if \lstinline|k >= size| the list is rotated by \lstinline|k % size|.
For example, both \lstinline|rotateRight(2)| and \lstinline|rotateRight(7)| turn \lstinline|[1, 2, 3, 4, 5]| into \lstinline|[4, 5, 1, 2, 3]|.
On an empty list, the method does nothing.

Your method must run in $O(N)$, and must not create any node nor change any \lstinline|data| field.

Hint: one way is to make the list circular, move \lstinline|head| and \lstinline|tail| forward together by the right number of steps, then cut the circle open again.

\begin{answerspace}{3in}
\ifanswers
\begin{lstlisting}
public void rotateRight(int k) {
    if (size == 0) return;
    tail.next = head;                    // make the list circular
    for (int i = 0; i < size - k % size; i++) {
        head = head.next;
        tail = tail.next;
    }
    tail.next = null;                    // cut the circle
}
\end{lstlisting}
\fi
\end{answerspace}
\end{question}
